% !TeX root = ../main.tex

\begin{abstract}

近年來，隨著量子計算的發展，擴展量子計算的規模和量子網路已成為研究重點之一。
本論文的目標是提出一類新型的量子電路，專門用於對量子態進行判別，並探討在去極化雜訊或更一般的量子雜訊存在下的設計選擇。

本研究利用強大的奈瑪克擴張定理（Naimark's dilation theorem）與等距映射合成
（isometry synthesis）演算法，建立一套用於有限個量子態判別的量子電路合成流程。
針對不同相位的相干態所設計的量子電路實驗顯示，當雜訊因子 $\lambda \leq 10^{-3}$ 時，仍具有
抗去極化雜訊的能力。

此外，本文亦探討多種正算子值測度（POVM）之半正定規劃（semidefinite programming）公式，用以
判別混合量子態。
期望本研究能為未來的量子電路最佳化與量子硬體設計提供新的啟發與方向。

\end{abstract}

\begin{abstract*}

In recent years, there have been efforts to scale up quantum computation with
quantum networks.
Our goal is to introduce a new class of quantum circuits that corresponds to
quantum state discrimination and to investigate the design choices when the
depolarizing noise or quantum noise in general is present.
In this thesis, I take advantage of the powerful Naimark's dilation theorem
and isometry synthesis algorithms and propose a quantum circuit synthesis
flow for discriminating a finite number of quantum states.
Quantum circuits designed to discriminate coherent states with different
phases show that they can resist depolarizing noise when the noise factor is
$\lambda \leq 10^{-3}$.
I also discuss different semidefinite programming formulations to calculate
POVMs for discriminating a set of mixed quantum states.
We hope that this work would inspire research on quantum circuit optimization
and quantum hardware design.

\end{abstract*}